\section{Sliding Window Inversions}
\label{sec:cses-sw-inversions}

\problemstatement{Given an array of $n$ integers and a window size $k$,
output the number of inversions in every contiguous window -- pairs of
positions $i<j$ within the window with $a_i>a_j$.}

\begin{patternbox}
Recomputing inversions per window is $O(nk\log k)$. Instead maintain a
running inversion count and a \textbf{Fenwick tree} over value ranks:
adding the new right element and removing the old left element each change
the count by a range query answerable in $O(\log n)$.
\end{patternbox}

\subsection*{Incremental inversion count}

Keep a Fenwick tree of the value counts currently in the window and a
running total $\textit{inv}$. When a new element $x$ enters on the
\emph{right}, it forms an inversion with every current element strictly
greater than $x$ (all of which sit to its left), so
$\textit{inv}\mathrel{+}=(\text{window size})-\#\{\le x\}$; then insert
$x$. When the leftmost element $y$ leaves, it was to the left of everything
else, so it formed an inversion with every current element strictly less
than $y$: $\textit{inv}\mathrel{-}=\#\{<y\}$; then remove $y$. Coordinate-
compress values so ranks fit the Fenwick tree.

\begin{ideabox}[Add on the Right, Remove on the Left, Query the Rank]
The window's geometry makes the updates clean: a right-entering element
pairs with larger existing elements; a left-leaving element pairs with
smaller ones. Both counts are Fenwick prefix queries, so each slide is
$O(\log n)$.
\end{ideabox}

\subsection*{Algorithm}

\begin{enumerate}
  \item Compress values to ranks $1\ldots m$. Build the first window,
        adding elements and accumulating inversions.
  \item Per slide: remove the leftmost ($\textit{inv}\mathrel{-}=$ count
        of smaller present), then add the new ($\textit{inv}
        \mathrel{+}=$ count of larger present).
  \item Output $\textit{inv}$ per window.
\end{enumerate}

\subsection*{Dry run}

\dryrun{$a=[2,1,3]$, $k=2$.}

\begin{itemize}
  \item Window $[2,1]$: $2>1$ is one inversion $\Rightarrow1$.
  \item Slide to $[1,3]$: remove $2$ (it had $1$ smaller element $\to
        \textit{inv}=0$), add $3$ (nothing greater present $\to$ stays
        $0$).
  \item Output $1,0$.
\end{itemize}

\subsection*{C++ implementation}

\begin{lstlisting}
#include <bits/stdc++.h>
using namespace std;

struct BIT {
    int n; vector<int> t;
    BIT(int n) : n(n), t(n + 1, 0) {}
    void update(int i, int v) { for (; i <= n; i += i & -i) t[i] += v; }
    int query(int i) { int s = 0; for (; i > 0; i -= i & -i) s += t[i]; return s; }
};

int main() {
    int n, k;
    cin >> n >> k;
    vector<int> a(n);
    for (int& x : a) cin >> x;

    vector<int> vals(a);
    sort(vals.begin(), vals.end());
    vals.erase(unique(vals.begin(), vals.end()), vals.end());
    auto rankOf = [&](int x) {
        return int(lower_bound(vals.begin(), vals.end(), x) - vals.begin()) + 1;
    };
    int m = vals.size();
    BIT bit(m);

    long long inv = 0;
    int inWin = 0;
    for (int i = 0; i < k; i++) {              // build first window
        int r = rankOf(a[i]);
        inv += inWin - bit.query(r);           // existing elements > a[i]
        bit.update(r, +1); inWin++;
    }
    cout << inv << " ";
    for (int i = k; i < n; i++) {
        int ry = rankOf(a[i - k]);             // remove leftmost
        inv -= bit.query(ry - 1);              // present elements < y
        bit.update(ry, -1); inWin--;

        int rx = rankOf(a[i]);                 // add new right
        inv += inWin - bit.query(rx);          // present elements > x
        bit.update(rx, +1); inWin++;
        cout << inv << " ";
    }
    cout << "\n";
    return 0;
}
\end{lstlisting}

\begin{complexitybox}
\textbf{Time Complexity.} $O(n\log n)$ -- each slide does a constant number
of Fenwick operations. \textbf{Space Complexity.} $O(n)$ for the Fenwick
tree and compression.
\end{complexitybox}

\begin{takeawaybox}
Sliding-window inversions become linear-log by maintaining a running count
plus a Fenwick tree of present values: the right-entering element pairs
with larger elements, the left-leaving with smaller. Recognising which side
an element joins/leaves fixes which prefix query to run.
\end{takeawaybox}
